% !TEX ROOT = thesis.tex
\chapter{Discussion}

\hspace{12.7mm}This chapter interprets and analyzes the research findings. It also provides a detailed discussion of the implications of the obtained results and the significance of the study. In addition, a comparison with the existing literature is included to highlight the contribution of this study to the field. The chapter is organized in multiple sections to discuss the results of the key methods in this research. 

\section{Automatic Ensemble Feature Selection}
The proposed automatic ensemble feature selection method was successful in its main goal. It consistently identified minimal and computationally efficient feature subsets for all five classification tasks. This is a key step for developing cost-effective models, which is a main objective of this research. This was achieved through the method's novel design, which works to find the number of features where adding more features provides a diminishing reduction in class overlap. This process ensures a balance between model simplicity and class separability. As a result, the method selected the smallest and most informative feature set. For instance, it reduced the feature count by around 82\% (from 68 to 12) for the normal cognition (\acrshort{NC}) vs. cognitive impairment (\acrshort{CI}) task and by approximately 85\% (from 71 to 11) for the cognitive impairment but not dementia (\acrshort{CIND}) vs. dementia task. This large reduction in features directly improved computational efficiency, making model training over 2.7 times faster than the baseline in both tasks. 

An important finding was related to the class overlap. The proposed method is designed to always select feature sets with the lowest F1 feature overlap scores. However, this reduction in class overlap did not directly result in the best predictive performance for all models. This was most clear in the future prediction tasks and the CI subtype classification task. In these tasks, the feature sets with the lowest class overlap produced a lower geometric mean (\acrshort{G-mean}) score with the random forest (\acrshort{RF}) model compared to the baseline and the Boruta method. This finding suggests a model-dependent effect. The F1 is a measure of linear separability based on Fisher's Discriminant Ratio. As a result, the method selected features that were highly effective for linear classifiers like logistic regression (\acrshort{LR}). In fact, for LR models, the minimal set selected by the proposed method performed equal to or better than the feature set selected by the Boruta method. 

In contrast, the non-linear RF model may have benefited from the larger, more complex feature sets of the baseline. Additionally, the Boruta method is a wrapper feature selection method based on an RF model. Therefore, it selects features that are optimized for an RF classifier. Because both the selection algorithm and the evaluation model are based on RF, it is logical that this method would produce a feature set that performs well for this specific model. On the other hand, the proposed method is model agnostic and optimized for a linear separability measure, which explains why it showed highly competitive performance with a linear model like LR. Therefore, the use of the proposed method aligns with the objective of this study since it contributes to improving the performance of linear models, which are generally more efficient and interpretable compared to non-linear models. 

Looking at the types of the selected features in different classification tasks, the analysis indicates a very strong and consistent pattern. For all four cognitive status classification and prediction tasks, the most dominant group of selected features was from the assessments of instrumental activities of daily living (\acrshort{IADL}) from the Functional Activities Questionnaire (\acrshort{FAQ}). This finding is logical. It aligns with the clinical definition of CI, which is primarily characterized by a decline in a person's ability to perform daily functional activities. In addition, similar studies in the literature also reported that IADL assessments are one of the top predictors of CI and dementia \cite{abbas2024Explainable, Ren2023Improving, Richter2025Cognitive}. \hl{This was also observed in prediction models from different countries, including LMICs} \cite{Cui2024Prediction, Ren2025Using}. Table~\ref{table:feature-details} in the appendix provides details about the selected features in all classification tasks.

A closer look at the IADL assessments shows that a constant set of features, such as BILLS, TAXES, SHOPPING, REMDATES, and TRAVEL, were selected in all four status tasks. \hl{The diagnostic value of these features is also supported by clinical evidence} \cite{Teng2010Utility}. Moreover, these functional features were consistently selected along with basic demographics (such as NACCAGE) and subjective memory complaint (MEMPROB). This finding has a high practical value. It suggests that in clinical practice, only a small set of data needs to be collected from the patient to run all the different models and provide a complete CI risk assessment. This makes the entire assessment workflow more cost-effective for application in clinical practice by reducing the cost and time required to gather the required data. 

For the more complex classification task of the CI subtypes (Alzheimer's disease (\acrshort{AD}) vs. non-AD), the feature set was dominated by the Clinical Dementia Rating (\acrshort{CDR}) scores, since all six CDR features were selected. The CDR scale is able to assess the nature and severity of the impairment in multiple cognitive domains. This finding indicates that the inclusion of CDR features is important to differentiate between AD and non-AD disorders as these classes cannot be effectively classified using demographic and functional assessment data. Similar studies also used CDR scores to differentiate AD from mild cognitive impairment \cite{abbas2024Explainable, Pang2023Predicting}. 

In summary, the proposed method successfully achieved its main goal of identifying effective and computationally efficient feature sets for all tasks. A key finding was that the method's model-agnostic design makes it particularly effective for linear models, which aligns with the research objective of developing cost-effective and interpretable models.

\section{Data Resampling}
The data complexity was analyzed in each of the tasks in this study to apply proper data resampling methods that could boost the performance of the machine learning models. The class imbalance ratio (\acrshort{IR}) was found to be relatively low in all tasks, ranging from 1.18 to 1.46. In contrast, class overlap was consistently high in the majority of tasks. The class instances were mostly overlapped at the global structural level, with the \acrshort{N1} scores ranging from approximately 0.30 in the current classification tasks to extreme values exceeding 0.50 in the CI subtype classification tasks. This suggests that utilizing only accessible data to predict cognitive impairment could result in structural intertwining of the classes in the feature space even if the class imbalance is minor. With demographic, depression scale, and cognitive assessment data, it is challenging to determine what cognitive status a subject might have in the future or to classify their CI subtypes as AD or non-AD. This is because many cases share similar values in these low discriminative data domains.

A comprehensive interpretation of various data resampling methods highlights the importance of handling class overlap, instead of only handling class imbalance, to obtain better classification performance using accessible data. In the NC vs. CI classification task, the baseline model struggled more in accurately classifying the majority class (CI) than the minority class (NC). This behavior indicates that due to the moderate class overlap, the model could not learn a distinct pattern of CI, even though there were more instances of CI than NC. As a result, despite how skewed the classes are, handling class overlap is necessary to obtain reasonable accuracy for both classes. Compared with other data resampling methods, the proposed method made the feature space easier to separate by removing noisy, confusing instances and strengthening the representation of the minority class in safe areas. This made the classes more distinct, resulting in slightly better recall in both classes.

The class overlap-guided data resampling method resulted in even better performance improvement in the future cognitive status prediction tasks. In these tasks, classes are highly intertwined because accessible data have limited discriminative power of future cognitive status. The proposed method employed the Edited Nearest Neighbor (\acrshort{ENN}) method, which analyzes the neighbors of all instances of both classes and removes instances surrounded by neighbors of the opposing class. After that, a suitable targeted oversampling method was applied to populate the sparse regions of the minority class. This process restructured the training data by eliminating confusing instances with similar feature values and systematically generating new instances that enhanced the definition of the decision boundaries. Consequently, the performance of the RF model improved significantly, especially in the future CIND vs. dementia prediction task, resulting in the best recall in both classes.

An opposite of class overlap guided data resampling is random resampling and Synthetic Minority Over-sampling Technique (\acrshort{SMOTE}). These standard techniques proved their ineffectiveness in the future cognitive status prediction tasks. For instance, Random Oversampling (\acrshort{ROS}) and SMOTE generated new instances in ways that worsened or did not resolve the class overlap (SMOTE increased the N1 score in the future NC vs. CI prediction task). ROS randomly duplicates minority class instances, including those in overlapped regions, whereas SMOTE creates synthetic instances by interpolating points within the complex boundary regions. Such oversampling approaches improved the RF model's recall of the minority class but considerably reduced its ability to correctly identify the majority class.

The Random Undersampling (\acrshort{RUS}) method also demonstrated similar deficiencies. While RUS could remove overlapped instances of the majority class, it could also remove informative instances in safe data regions, resulting in more complex data. For instance, RUS resulted in the worst dementia prediction accuracy of 0.723 compared to 0.77 of the baseline and 0.776 of the proposed method. This trade-off could be undesirable in certain real-world clinical scenarios, as it increases the risk of missing actual dementia cases.

By analyzing the relationship between class overlap measures and the improvement in the RF performance, it can be realized that handling class overlap is particularly valuable in classification problems characterized by high structural complexity. This relation was mainly observed in the AD vs. non-AD classification task, which presented the highest degree of structural class overlap. In this task, the proposed method achieved the highest recall for both the AD and non-AD classes simultaneously.

\newpage
Data cleaning by ENN is a major factor in the success of the proposed method. While ENN helped the model focus on clear patterns by eliminating confusing and probably noisy instances, it could have possibly removed informative instances that represent valid characteristics of both classes. However, in general, this method showed potential in improving the quality of the training data of the National Alzheimer's Coordinating Center (\acrshort{NACC}) Uniform Data Set (\acrshort{UDS}), where subjects might have been misdiagnosed or experienced fluctuating cognitive status.

It is important to note that in certain tasks, such as the prediction of future NC vs. CI, the overall improvement in the G-mean was marginal, even though the class overlap was successfully resolved. This slight increase could be due to the inherent informational limitation of the data used in this study. Although the proposed method maximized the learnability of the data by minimizing class overlap, it was unable to generate information that was absent from the feature set. Further enhancement of the model's accuracy in predicting future cognitive status or classifying CI subtype might require more data that bring fresh and different information beyond demographics and basic cognitive assessments.

In summary, the results demonstrated that poor performance attributed to low-cost, accessible data is largely due to structural class overlap. It is possible to resolve this overlap and produce accurate and clinically viable predictive models of CI without the use of expensive biomarkers by applying the designed data resampling method. This supports the research objective of developing accessible CI screening models in resource-constrained settings.

\section{Model Fine-Tuning and Selection}

Following the data preprocessing phase, five different machine learning algorithms were optimized to select the best set of hyperparameters that can provide a balance between predictive accuracy and computational efficiency. The selection was automated using the novel cost-effective model selection algorithm (Algorithm \ref{alg:model-selection}). This section discusses the key findings of this process.

\subsection{Efficacy of the Cost-Effective Model Selection Algorithm}
The application of the proposed selection algorithm resulted in significant improvements in computational efficiency across various models. The algorithm allowed for the selection of cost-effective models, which increased the prediction speed without causing a statistically significant decrease in classification performance. For instance, in the future CIND vs. dementia prediction task, the prediction throughput of the selected cost-effective k-Nearest Neighbors (\acrshort{KNN}) model was increased by 3.15 compared to a model with the best performance. Similarly, in the future NC vs. CI prediction task, the training speed of the Light Gradient Boosting Machine (\acrshort{LightGBM}) model was improved by 3.59 times. These gains were achieved while keeping the performance reduction strictly within 1\% of the best obtained performance. Section~\ref{sec:lr-fine-tuning-apdx} in the appendix demonstrates how the proposed model selection algorithm was able to select the most cost-effective LR models in all classification tasks.

These results suggest that the benchmark models, which were identified by the algorithm based on the best classification accuracy, were often over-parameterized. The selection algorithm successfully selected hyperparameters that are able to learn the essential patterns in the training data without unnecessary computational complexity. For instance, the selection algorithm continuously selected shallower trees for ensemble models and simplified distance metrics for KNN models.

Beyond the scope of this study, the proposed cost-effective model selection algorithm offers more valuable advantages in terms of its practical use. The structured three-phase approach of the algorithm streamlines and automates the complexity of the model selection process and promotes model selection criteria beyond classification performance. This structured design allows technical and non-technical users to easily understand how the algorithm works and produce the intended output. Moreover, the algorithm is designed with multiple parameters to define different performance and efficiency requirements. This indicates that the algorithm is highly adaptable to various industrial and academic applications. Furthermore, the algorithm can be easily integrated into existing machine learning pipelines, which strengthens its practical utility. Section~\ref{sec:selection-algorithm-usecases-apdx} in the appendix discusses how the proposed cost-effective model selection algorithm can be utilized in various use cases.

\subsection{Final Model Selection}
Based on the comparative analysis, LR was selected as the final model for all classification tasks in this study. It was observed that the LR model achieved performance that was statistically comparable to complex non-linear models, such as RF, LightGBM, and Support Vector Machines (\acrshort{SVM}). For instance, in the NC vs. CI classification task, the LR model achieved an F1-score of 0.848, which was very close to that of the RF model (0.849) and better than the scores of the other models.

This finding contrasts with the results reported in most similar studies in the literature, where linear models like LR perform worse than complex and ensemble models. There is a common assumption in the domain of medical predictive modeling that deep learning and ensemble methods often outperform simple and linear models \cite{Christodoulou2019systematic}. As a result of this assumption, multiple CI and dementia prediction studies did not involve the evaluation of linear models and ignored the computing efficiency value that these simple models may provide. While this practice could be valid for unstructured data like Magnetic Resonance Imaging (\acrshort{MRI}) scans, it is not necessarily the best practice in predictive problems based on structured clinical data, especially when the objective is to develop generalizable and accessible models. 

The comparable accuracy of the LR models in this study can be attributed to the use of feature selection and data resampling methods that were designed to reduce class overlap. These methods systematically eliminated irrelevant features and removed noisy data, making the feature space clearly structured and more linearly separable. \hl{In such a feature space, a simple and low-complexity model like LR becomes theoretically favorable. The high capacity of complex ensembles is no longer needed to capture a simpler decision boundary, and this redundant capacity could increase the model variance and the risk of overfitting. In contrast, a low-complexity model is expected to perform better on unseen data.} To support this finding, a systematic review \cite{Christodoulou2019systematic} concluded that machine learning models do not provide a performance benefit over LR for clinical prediction models. Similarly, a more recent study \cite{Wu2023Logistic} \hl{reported that LR achieved a performance comparable to that of complex machine learning models in predicting cognitive impairment using low-dimensional clinical data.}

In addition to comparative performance, there are other advantages that encourage the selection of the LR model. First, the LR model is computationally efficient because it is designed based on a simple mathematical formula to map the input to the intended prediction. This calculation requires minimal time and computational resources compared to other models that need to perform repeated complex calculations. The efficiency of the LR model is demonstrated in the results by achieving training and prediction speeds that significantly outperformed other models in all classification tasks. 

The second advantage is that the simple structure of the decision process in the LR model makes it transparent and inherently interpretable \cite{Rudin2019Stop}. The relationship between the risk factors (features) and the predicted cognitive status can be clearly explained by understanding the LR model's coefficients. In contrast, complex and ensemble algorithms are considered black-box systems that require a post-hoc explanation of their decision. This could be considered a bad practice and could generate misleading explanations that incorrectly reflect the model decision process \cite{Rudin2019Stop}.

Finally, the selection of the LR model could enhance generalizability to new unseen data. Complex models could memorize noise in the training data due to their high learning capacity. To avoid this risk of overfitting, additional development effort and cost might be required to collect more training data or apply specialized techniques, such as regularization. The generalizability of the LR model is discussed in the next section.

It can be concluded that the results of model fine-tuning and selection support the objectives of this study to develop cost-effective models for CI classification and future risk prediction. The analysis of the results indicates that while the selected LR models achieve high accuracy, they also offer outstanding computational efficiency, inherited interpretability, and effective generalizability. This implies that the solution can be developed at reduced cost and can be deployed on low-resource hardware, such as basic smartphones or tablets, without the need for expensive cloud infrastructure. As a result, this could improve the accessibility of the solution and boost the security and privacy of personal data. In addition, the inherited interpretability of the LR model is essential to build trust with clinicians and patients in real-world healthcare settings.

\section{External Validation}
After selecting the most cost-effective model for each classification task, the models were evaluated on two independent datasets. \hl{Overall, the models presented high discriminative performance and clinical utility in the evaluated datasets, suggesting their potential suitability for cognitive status screening in clinical settings similar to those studied.} The results are discussed in more detail in the following subsections.

\subsection{Assessment of Classification Generalizability}
The evaluation of the selected LR models in the external validation showed that the models achieved strong performance across most of the classification tasks. There was no significant difference between the area Under the Curve (\acrshort{AUC}) scores in the internal cross-validation (CV) and the scores of the two external validations in the current cognitive status classification tasks. More interesting results were found in the future prediction tasks, as the models' performance in external validation was higher than their performance in the internal CV. For instance, the models achieved AUC scores ranging from 0.78 to 0.80 in external validation, compared to approximately 0.77 in the internal CV.

This finding challenges the common results reported by similar studies in the literature, where the performance of the model usually decreases when applied to unseen external data. The positive generalizability performance obtained could be the result of handling class overlap through the proposed systematic data resampling method. By removing confusing or noisy samples from the training dataset, the model learned a clear decision boundary that generalized well to external datasets.

\newpage
In addition, it was noticed that the models performed better in the Alzheimer's Disease Neuroimaging Initiative (\acrshort{ADNI}) external validation than in the NACC cross-center validation for classifying the current cognitive status. This difference is likely due to the characteristics of the data sources. Data samples in the ADNI dataset might have a clearer clinical profile compared to those in the NACC UDS, which may include some noisy samples due to misdiagnosis.

Despite the excellent generalizability to external data, a low performance was observed in the CI subtype classification (AD vs. non-AD). The model achieved an AUC of 0.670 in the NACC cross-center validation, which is considered below the clinically accepted performance for screening models. The reason for this low performance could be due to the significant feature overlap between AD and other CI subtypes when only accessible features like demographics, cognitive assessments, and function assessments are used. Additionally, these features were not mainly used to diagnose the primary etiology of CI in the NACC UDS. Therefore, these accessible features are expected to have low discriminative power, which confirms the hypothesis of this study. To accurately classify the primary cause of CI, biomarkers, such as MRI or genetic testing, are required. However, these biomarkers are expensive and not widely available in primary care. This limitation is supported by existing literature, which highlights that clinical data that are accessible in primary care could be sufficient for accurate screening of cognitive status, but might be insufficient to accurately identify the subtypes of CI \cite{Yineng2019Machine, Zhu2020Machine}.

In general, the results confirm that accessible data are sufficient for broad screening and risk assessment, but not sufficient to determine the specific cause of CI. Therefore, the proposed models are best suited for primary care settings to identify those with impairment and the severity of their cognitive impairment, but the accuracy of identifying the subtype of CI could be limited; therefore, the model output should not be used as the main source of information to make final clinical decisions.

\subsection{Assessments of Model Calibration and Clinical Utility}
In addition to assessing the performance of the models in external datasets, it is important to ensure that the models produce predicted probabilities that accurately reflect the actual risk. The NACC cross-center validation results indicated that the uncalibrated models consistently made overconfident predictions, with calibration slopes much lower than 1.0, except for the CIND vs. dementia classification task. This behavior could be a result of the data resampling strategy used to handle class overlap. According to a recent study \cite{Piccininni2024Understanding}, models might produce extreme probabilities when trained on balanced datasets that were resampled to improve model performance. This is because, in the cleaned training data, the models were not exposed to the level of uncertainty present in real-world data. This confirms that the calibration applied in this study is a necessary step to ensure that the prediction is statistically meaningful. Calibration is crucial for machine learning models' reliability and effectiveness in clinical use; nevertheless, it was not reported in the majority of the studies in the literature.

The calibration of the model in the NACC cross-center validation resulted in Brier scores below 0.15 in all four cognitive status classification tasks. These scores indicate that the models are both accurate and well calibrated. Moreover, the excellent calibration was also noticed by the Expected Calibration Error (\acrshort{ECE}) around 0.05 in most of the classification tasks. The models estimate the targeted risk with probabilities that align with the actual risk, with an average deviation of only 5\%. This result was achieved by calibrating the models' probabilities using beta calibration without the need to retrain the models. 

The outstanding model calibration result in the NACC cross-center validation indicates that when diagnosis criteria and data collection protocols are standardized, the calibration can generalize effectively across different geographical centers. Despite the excellent calibration results in the NACC cross-center testing data, the ECE in the CIND vs. dementia prediction task was higher than in the other classification tasks. Although the ECE of 0.074 is still considered acceptable, this relative increase in the calibration error could be due to the unclear boundaries in the feature space when it comes to classifying future CIND cases from early dementia cases using low-cost clinical and demographic data.

On the other hand, the models presented poor calibration in the ADNI external validation, although they achieved high performance in most of the classification tasks in the same external validation. This issue is known as "calibration drift" and occurs when the prevalence of the disease in the new setting differs significantly from the training data. This finding aligns with the argument of \cite{VanCalster2019Calibration} that a model can be accurate in ranking patients but unsafe in estimating risk if it is not recalibrated. Therefore, recalibration is a mandatory step before deploying these models in new populations with different frequencies of the targeted condition. Otherwise, the probabilities of the risk should not be interpreted as precise estimations when the model is used in different populations.

The clinical utility of the models were further evaluated using Decision Curve Analysis (\acrshort{DCA}). The DCA results of most classification tasks demonstrated that the models provided a better net benefit than default strategies like "treat all" or "treat none" across a wide range of probability thresholds, usually between 0.11 and 0.99. \hl{These results mean that the probability threshold can be tuned to the clinical goal of a specific subgroup. This clarifies the practical meaning of the sensitivity and specificity trade-off. One example is a screening subgroup of older adults in low-resource primary care. For this subgroup, a low threshold is preferred to keep the model at high sensitivity because a missed case of cognitive impairment is more harmful than a false alarm that a simple follow-up can clear in a later stage. Another example is a confirmation subgroup of patients who already show symptoms before referral. For this subgroup, a high threshold is preferred to keep the model at high specificity. In this case, patients are referred to costly assessments, such as brain scans, only if they have a high probability of risk. Consequently, the wide probability threshold range allows the models to serve different clinical diagnostic workflows and to support cost-effective screening.}

In conclusion, the external validation results of the cost-effective models confirm the achievement of the research objectives. The models demonstrated robust performance and high clinical utility for cost-effective screening and future risk prediction. The results also confirm the research hypothesis that accurately classifying CI subtypes using accessible data is a challenging task. Additionally, it has been noted that local recalibration is required when deploying the models in environments with different disease prevalence.

\section{Model Explainability}
A comprehensive evaluation of explainability at multiple levels, including global, local, and counterfactual explainability, was conducted on the final models. The purpose of this method is to ensure clinical interpretability of the model predictions after verifying the models' strong generalizable performance and clinical utility in the external validation phase. The overall results indicate that the developed models rely on clinically relevant features, and they can be used to generate insightful and actionable explanations at the individual level for CI risk management and prevention. 

It should be noted that the LR model was selected for all classification tasks, which is inherently interpretable. The application of post-hoc explainability methods such as the SHapley Additive exPlanations (\acrshort{SHAP}) is not required to provide global and local explanations of the LR model's prediction. The coefficients of the model can be directly used instead. However, SHAP was used in this study to provide a consistent report of the model explanation, which is commonly used in the literature.

\subsection{Assessment of Global Explanation}
There was a high degree of similarity found between the NACC cross-center validation and the ADNI external validation in the global feature importance rankings. The Spearman’s rank correlation coefficient was more than 0.83 across all classification tasks. This consistency could indicate that the models are not overfit to training data and are capable of learning reliable patterns. \hl{This further supports the evidence that the proposed framework, including the systematic feature selection and data resampling methods, contributes to the development of generalizable and reliable cost-effective models, although broader validation across more diverse clinical populations remains necessary.}

The study of the direction of feature influence based on SHAP (SHapley Additive exPlanations) summary plots showed that a higher dependence on performing functional activities, as well as active depression (\acrshort{GDS}) and behavioral symptoms, are consistently linked to a higher risk of cognitive impairment and dementia. Regarding demographic and physical measures, older age and male gender are found to relatively increase the estimated risk, whereas the low body mass index (\acrshort{BMI}) measures are associated with the early risk of dementia. Section~\ref{sec:xai-apdx} in the appendix includes the SHAP summary plots for all classification tasks showing the direction of the impact of the features on the model prediction.

Despite the strong similarity in global feature importance between the two external validation cohorts, notable differences were observed. For instance, in the Future NC vs. CI prediction task, irritability (IRR) was ranked significantly higher in the ADNI external validation compared to the NACC cross-center validation. This difference is likely attributed to cohort heterogeneity, as the ADNI cohort may include larger groups with behavioral symptoms compared to the NACC cohort.

Regarding the differences among classification tasks, it was observed that features related to memory assessment, such as REMDATES and MEMPROB, were found to be the most important indicators of early impairment in the NC vs. CI classification and prediction tasks. In comparison, complex executive functions, such as paying bills (BILLS), were the top predictors to identify the severity of CI in the CIND vs. dementia classification and prediction tasks. This finding aligns with the hierarchy of functional loss described in the literature, where complex instrumental activities are lost after basic memory functions \cite{Njegovan2001Hierarchy}. Furthermore, another variation was observed between the current classification and the future prediction tasks. While current diagnosis relies on objective deficits, such as REMDATES, future prediction is more influenced by subjective memory complaints (MEMPROB). This was also noticed by a study highlighting that subjective complaints appear in the early phase of CI before objective clinical measurement \cite{Jessen2014Conceptual}.

In addition to these differences, functional assessment (\acrshort{FAQ}) features were the most important features in most of the classification tasks. This could be because of their clear and strong link to the subject's cognitive status compared to other behavioral and demographic features. Finally, the analysis of feature rankings pointed out that demographic features, such as age and sex, were commonly ranked lower than functional and behavioral symptoms in most of the classification tasks. However, age was found to be more important than some CDR features in the AD vs. non-AD subtype classification.

\subsection{Assessment of Local Explanation}
The study of local explanations showed that in all classification tasks, different people could be given the same predicted probability of having CI or dementia, even though their risk profiles are very different. For instance, a subject with difficulty performing financial tasks could be assigned the same probability of CI as another subject who is older and suffers from memory loss. A recent study supports this observation, highlighting that dementia is a diverse condition that could be caused by various symptoms \cite{Vogel2021Four}. As a result, local explanation is important in identifying specific risk factors on an individual basis for effective communication of the risk and informed clinical decision-making.

Moreover, it can be concluded that global feature importance cannot be used to interpret individual predictions. Features that were ranked low in the global feature importance, such as age, were found to be one of the primary drivers of prediction for specific individuals according to the analysis of SHAP waterfall plots. This occurs because global importance represents an average across the population, whereas local importance captures specific contributions for a single instance. Therefore, local explanation is a crucial component in the development of CI screening tools \cite{Hakkoum2024Global}.

\subsection{Assessment of Counterfactual Explanation}
The counterfactual analysis demonstrated that the Diverse Counterfactual Explanations (\acrshort{DiCE}) method was successful in generating valid counterfactual explanations for the external ADNI dataset. DiCE generated counterfactual examples for more than 86\% of the instances in all cognitive status classification tasks, with a success rate of up to 100\% in the NC vs. CI prediction task. This success could imply that the models, which were trained on NACC data, are able to find feasible alternative scenarios (counterfactuals) for subjects in a completely different cohort (\acrshort{ADNI}). \hl{This further suggests that the developed models may generalize across different datasets and feature spaces within similar clinical settings.}

It should be noted that DiCE encountered challenges in explaining advanced cognitive impairment stages, specifically in the CIND vs. dementia prediction task. The \acrshort{KD-Tree} (K-dimensional tree) search algorithm did not generate counterfactuals using real-world data points. This difficulty could be due to the sparsity of data and the high overlap between CIND and dementia, especially when using accessible data to predict the future risk of dementia. KD-Tree had to spend substantial time looking for neighbors in such an ambiguous data domain before it failed to find any valid neighbors for about half of the instances. This limitation is consistent with literature stating that counterfactual analysis in heterogeneous diseases, such as dementia, is often constrained by the lack of complete health records \cite{Kim2022Counterfactual}. Therefore, it was required to use the random search method instead, although it generates synthetic examples that might not reflect real-world data. 

From the analysis of counterfactual explanation, it was observed that early prevention of CI risk is easier and more feasible compared to reversing current impairment. According to the results of future prediction tasks, preventing most of the predicted future risks required modifying only 1 to 3 features. In contrast, reversing a current diagnosis required changing 5 to 6 features on average, according to the results of current classification tasks. To technically explain this finding, subjects with future risk are located closer to the decision boundary and are easier to shift back to the normal class region compared to those with established impairment. It is agreed that risk modification in the early phase is feasible, whereas reversing established pathology is currently impossible \cite{Livingston2020Dementia}. Therefore, this observation confirms that early intervention efforts before the first onset of CI are crucial to completely prevent the risk, which aligns with the objective of this study.

Furthermore, it was observed that the feature importance generated by DiCE consistently ranks the Instrumental Activities of Daily Living (IADL), such as financial management (BILLS and TAXES), as the most important features to reverse the risk. This agrees with studies identifying financial capacity as a critical measure of conversion to dementia \cite{Jekel2015Mild}. According to the SHAP feature importance rank, these features have a high discriminative power, which could explain why DiCE finds them to be the most efficient features to change in order to get the most significant reduction in the predicted risk probability with the minimal required modifications.

Although DiCE mathematically succeeded in generating counterfactual examples by overrelying on IADL features, there is a concern about whether improvement in IADL is always possible and how clinicians and patients benefit from such explanations. While improvements in IADL might not be possible, especially when the underlying neurodegenerative damage is irreversible, the suggestion to modify IADL can be interpreted as a guide to find practical workarounds to help the patient manage their daily life instead of trying to cure the biological damage. Hence, clinicians should utilize these findings to implement targeted environmental modifications, such as prescribing occupational therapy or financial guardianship, to stabilize the subject's functional status. Moreover, patients and their caregivers should be advised to work on reducing the burden caused by these specific tasks. For instance, automating or delegating financial management tasks can help mitigate the impact of the impairment and maintain quality of life.

In conclusion, the explainability analysis successfully achieved the study's objective to provide a clear and insightful interpretation of the model's output that supports preventing the CI risk. With the combination of global, local, and counterfactual explanations, the proposed method is not limited to explaining why a prediction is made but also explains how it applies to the individual subject and what actions might be possible to mitigate the risk.

\section{Performance Comparison with Existing Research}\label{sec:performance-comparison}
This section discusses the model performance results obtained in this study in comparison to the performance achieved by similar studies in the literature. The studies included in this comparative analysis are selected because they rely mainly on low-cost and accessible data to predict current or future cognitive conditions. This comparison is limited by the difference in classification strategy between this study and the majority of the previous research. This study adopts a hierarchical classification approach by classifying cognitive impairment first (NC vs. CI) before further classifying the severity of cognitive impairment to CIND vs. dementia. On the other hand, similar studies that use accessible data mainly classify normal cases from dementia cases. In addition, none of the compared studies had a CIND class in their classification tasks; instead, the severity of cognitive impairment is commonly classified as mild cognitive impairment (\acrshort{MCI}) or dementia. Refer to Table~\ref{table:ci-review} and Table~\ref{table:et-review} for more details of the referenced studies.

The performance of the LR models in this study is compared with other models in the literature as presented in Table~\ref{table:compare-classify}. In regard to studies that classify NC vs. CI, \citeA{Ren2023Improving} also utilized hierarchical classification and the NACC UDS to classify NC vs. CI. The higher accuracy of 0.889 was reported by \cite{Ren2023Improving} compared to the accuracy obtained in this study (0.833). This is because their random forest model was trained on data from the Mini-Mental State Examination (\acrshort{MMSE}), which are direct and strong indicators of CI. Moreover, \citeA{Ren2023Improving} only reported the accuracy of the model, which could be misleading, as accuracy is not a reliable measure when classes are imbalanced.

\begin{table}[!ht]
    \centering
    \caption{Comparison of model performance with the literature for the classification of current cognitive status}
    \begin{tabular}{|M{0.2\textwidth}|M{0.3\textwidth}|M{0.1\textwidth}|M{4cm}|}
    \hline
        Study & Target Classes & Model & Performance \\ \hline
        \textbf{This Study} & \textbf{NC (21,931) vs. CI (31,584)} & \textbf{LR} & \textbf{AUC = 0.844} \newline \textbf{Accuracy = 0.833} \\ \hline
        \cite{Li2023Machine} & NC (1,842) vs. CI (384) & \acrshort{GLM} & AUC = 0.779 \\ \hline
        \cite{Ren2023Improving} & NC (1,129) vs. CI (789) & RF & Accuracy = 0.889 \\ \hline
        \cite{Yue2023Explainable} & NC (2,004) vs. MCI (556) & AdaBoost & AUC = 0.957 \\ \hline
        \cite{Cui2024Prediction} & NC (1,008) vs. CI (1,130) & LR & AUC=0.875 \\ \hline
        \cite{Richter2025Cognitive} & NC (101) vs. CI (182) & SVM & AUC = 0.850 \\ \hline
        \cite{Ye2025Predicting} & NC vs. CI (n = 13,906) & GBC & AUC = 0.832 \\ \hline
        \textbf{This Study} & \textbf{CIND (14,432) vs. Dementia (17,164)} & \textbf{LR} & \textbf{AUC = 0.828} \newline \textbf{Accuracy = 0.828} \\ \hline
        \cite{Bucholc2023Hybrid} & MCI (656) vs. Dementia (112) & RF & AUC = 0.692 \\ \hline
        \cite{Ren2023Improving} & MCI (347) vs. AD (442) & RF & Accuracy = 0.810 \\ \hline
        \cite{abbas2024Explainable} & MCI (10,043) vs. AD (15,687) & SVM & Accuracy = 0.866 \\ \hline
        \multicolumn{4}{p{\textwidth}}{\footnotesize The reported performance is based on the internal validation results. The accuracy is used if the study did not report AUC.\newline
        NC: normal cognition, CI: cognitive impairment, CIND: cognitive impairment but not dementia, MCI: mild cognitive impairment, LR: logistic regression, GLM: generalized linear models, RF: random forest, SVM: support vector machines, GBC: gradient boosting classifier, AdaBoost: adaptive boosting, AUC: area under the curve, n: total sample size when the class split was not reported.} \\
    \end{tabular}
    \label{table:compare-classify}
\end{table}

Other studies, such as \cite{Cui2024Prediction} and \cite{Richter2025Cognitive}, also achieved better performance with AUC of 0.875 and 0.850, respectively. However, the models in both studies had a precision of 0.808 and 0.690, respectively, significantly lower than the precision of the proposed LR model (0.922). This indicates that while these studies achieved high discriminative performance, their ability to minimize false positives was limited. On the other hand, \citeA{Li2023Machine} reported the lowest AUC of 0.779, which could be due to the ineffective handling of the major class imbalance. \hl{In a recent study, accessible data from the China Health and Retirement Longitudinal Study (CHARLS) were used by} \citeA{Ye2025Predicting} \hl{to classify the current cognitive status, where a gradient boosting model reached an AUC of 0.832, close to the performance of the proposed LR model. A much higher accuracy was reported by} \citeA{Yue2023Explainable} \hl{for separating normal cognition from MCI; however, this result was obtained on a smaller cohort with strong oversampling and without external validation.}

In regard to the classification of CI severity, the LR model in this study performed better than the random forest models used in \cite{Bucholc2023Hybrid, Ren2023Improving}, even though those models were trained on data from cognitive tests and the MMSE. In contrast, \citeA{abbas2024Explainable} reported an accuracy of 0.866 in the internal validation using NACC UDS, outperforming the accuracy of the proposed LR model. The higher accuracy achieved by the SVM model in \cite{abbas2024Explainable} could be attributed to the use of CDR scores, which were reported as the most important features in their study. As a result, the SVM model also achieved higher accuracy in the ADNI external validation.

A comparison of the performance of the proposed LR prediction models is provided in Table~\ref{table:compare-predict}. All studies classify NC vs. dementia or NC vs. Alzheimer's disease and related dementias (\acrshort{ADRD}), excluding cases of MCI. This could be one of the reasons why few studies \cite{Akter2025Using, Georgiev2024Predicting, Li2023Early} achieved a higher AUC, in addition to the use of ensemble models. Moreover, \citeA{Georgiev2024Predicting} utilized a frailty index to train their model, which is a cumulative measure of deficits that consistently ranked as the top feature for predicting incident dementia. The other study \cite{Li2023Early} reported a low precision around 0.50, indicating that their model produces a high rate of false positives. On the other hand, the LR model developed in this study outperformed the random forest models of \cite{Ben2020Predicting, Park2020Machine}. \hl{Other studies that depend on accessible data reported a comparable or lower performance.} \citeA{Twait2023Dementia} \hl{obtained an AUC of 0.710 using only clinically accessible variables. The balanced random forest of} \citeA{Ren2025Using} \hl{reached a high accuracy of 0.885 but a low sensitivity of 0.580, which reflects a weak detection of the positive cases because of a high class imbalance. The effect of class imbalance was also observed in} \citeA{Oh2024Multivariable}\hl{, where only 70 dementia cases were available against 4,905 normal cases.}

\begin{table}[!ht]
    \centering
    \caption{Comparison of model performance with the literature for the prediction of future cognitive impairment}
    \begin{tabular}{|M{0.2\textwidth}|M{0.3\textwidth}|M{0.15\textwidth}|M{3.5cm}|}
    \hline
        Study & Target Classes & Model & Performance \\ \hline
        \textbf{This Study} & \textbf{NC (6,611) vs. CI (5,734)} & \textbf{LR} & \textbf{AUC = 0.772 \newline Accuracy = 0.776} \\ \hline
        \cite{Ben2020Predicting} & NC (11,558) vs. Dementia (2,159) & RF & Accuracy = 0.629 \\ \hline
        \cite{Park2020Machine} & NC (614) vs. AD (614) & RF & AUC = 0.725 \\ \hline
        \cite{Li2023Early} & NC (1,038,643) vs. ADRD (23,835) & GBT & AUC = 0.853 \\ \hline
        \cite{Twait2023Dementia} & NC (3,901) vs. Dementia (892) & ENR & AUC = 0.710 \\ \hline
        \cite{Oh2024Multivariable} & NC (4,905) vs. Dementia (70) & XGBoost & AUC = 0.819 \\ \hline
        \cite{Georgiev2024Predicting} & NC (132,670) vs. Dementia (11,143) & XGBoost & AUC = 0.888 \\ \hline
        \cite{Akter2025Using} & NC (119,723) vs. ADRD (4,012) & GBT & AUC = 0.833 \\ \hline
        \cite{Ren2025Using} & NC vs. CI (n = 2,688) & BRF & AUC = 0.809 \newline Accuracy = 0.885 \\ \hline
        \textbf{This Study} & \textbf{CIND (2,335) vs. Dementia (3,398)} & \textbf{LR} & \textbf{AUC = 0.759 \newline Accuracy = 0.750} \\ \hline
        \cite{Pang2023Predicting} & MCI (176) vs. AD (362) & RF & AUC = 0.805 \\ \hline
        \cite{abbas2024Explainable} & MCI (922) vs. AD (1,468) & RF & Accuracy = 0.767 \\ \hline
        \multicolumn{4}{p{\textwidth}}{\footnotesize The reported performance is based on the internal validation results. The accuracy is used if the study did not report AUC.\newline
        AD: Alzheimer's disease), ADRD: Alzheimer's disease and related dementias, NC: normal cognition, CI: cognitive impairment, CIND: cognitive impairment but not dementia, MCI: mild cognitive impairment, GBT: gradient boosting trees, LR: logistic regression, RF: random forest, SVM: support vector machines, XGBoost: eXtreme gradient boosting, ENR: elastic net regression, BRF: balanced random forest, AUC: area under the curve, n: total sample size when the class split was not reported.} \\
    \end{tabular}
    \label{table:compare-predict}
\end{table}

Regarding the prediction of future CI severity, the proposed LR model scored slightly higher accuracy than the random forest model in \cite{abbas2024Explainable}, despite the use of CDR to train the model. The performance of the random forest model was significantly reduced in the ADNI external validation compared to the LR model in this study, which maintained the same level of prediction accuracy. This demonstrates the potential of the proposed framework in developing generalizable models. The random forest model in \cite{Pang2023Predicting}, which was also trained using NACC UDS, achieved a higher AUC of 0.805 because it utilizes highly predictive features, such as the mode of onset of cognitive symptoms and CDR.

The performance comparison for the classification of the CI subtypes is presented in Table~\ref{table:compare-etology-classification}. Studies that use accessible data in this domain are limited. \citeA{Gurevich2017Neuropsychological} classified AD from other etiologies of dementia using demographic data and conventional clinical neuropsychological testing. The SVM model achieved an accuracy of 0.680, which is less than the accuracy achieved by the LR model developed in this study (0.700). Other studies focus on classifying AD from other specific types of CI such as vascular dementia (\acrshort{VaD}) and dementia with Lewy bodies (\acrshort{DLB}). \citeA{Davoudi2021Classifying} used kinematic, time‑based, and visuospatial features extracted from the digital Clock Drawing Test (dCDT) to classify AD patients from those with VaD with an accuracy of 0.768. \citeA{Yamada2022Characteristics} also utilized features extracted from drawing exercises and a RF model to classify AD vs. DLB with an accuracy of 0.738. Although these two studies reported better accuracy compared to this study, they used complex models (RF) and specialized tests that might not be available in primary care units. In contrast, \citeA{Sharma2024Leveraging} differentiated AD and VaD patients from DLB patients using accessible comorbidity and demographic features. They achieved a accuracy of 0.700 using an XGBoost model. \hl{In a multicentric study across Latin American countries,} \citeA{Maito2022Classification} \hl{classified AD from frontotemporal dementia (FTD) with an accuracy of 0.932 using routine clinical and cognitive measures. Although this accuracy is higher than the one obtained in this study, the random forest model was not validated on external data.} All the compared studies in Table~\ref{table:compare-etology-classification} did not validate their models in external data compared to this study, which validated the LR model in a cross-center dataset with acceptable performance. In addition, most of the compared studies used small datasets, which limits the reliability and the generalizability of the obtained results.

\begin{table}[!ht]
    \centering
    \caption{Comparison of model performance with the literature for the cognitive impairment subtype classification}
    \begin{tabular}{|M{0.2\textwidth}|M{0.3\textwidth}|M{0.15\textwidth}|M{3.5cm}|}
    \hline
        Study & Target Classes & Model & Performance \\ \hline
        \textbf{This Study} & \textbf{AD (18,201) vs. non-AD (13,441)} & \textbf{LR} & \textbf{Accuracy = 0.700} \\ \hline
        \cite{Gurevich2017Neuropsychological} & AD (70) vs. non-AD (88) & SVM & Accuracy = 0.680 \\ \hline
        \cite{Davoudi2021Classifying} & AD (29) vs. VaD (27) & RF & Accuracy = 0.768 \\ \hline
        \cite{Yamada2022Characteristics} & AD (47) vs. DLB (27) & RF & Accuracy = 0.738 \\ \hline
        \cite{Maito2022Classification} & AD (904) vs. FTD (282) & RF & Accuracy = 0.932 \newline AUC = 0.965 \\ \hline
        \cite{Sharma2024Leveraging} & AD/VaD (33,318) vs. DLB (1,657) & XGBoost & Accuracy = 0.700 \\ \hline
        \multicolumn{4}{p{\textwidth}}{\footnotesize The reported performance is based on the internal validation results.\newline
        AD: Alzheimer's disease, VaD: vascular dementia, DLB: dementia with lewy bodies, FTD: frontotemporal dementia, LR: logistic regression, SVM: support vector machines, RF: random forest, XGBoost: eXtreme gradient boosting, AUC: area under the curve.} \\
    \end{tabular}
    \label{table:compare-etology-classification}
\end{table}

Comparing the overall quality of the study design, our study is designed based on robust data preprocessing, model selection, and model validation methods that support the development of cost-effective and reliable models. \hl{The comparison across the main dimensions of methodological quality is summarized in} Table~\ref{table:compare-summary} \hl{in the appendix.} None of the comparable studies measured or handled the potential class overlap, despite the major class imbalance observed in some studies. \hl{Class imbalance was commonly addressed only by resampling or class weighting, while complex and ensemble models were used to boost the predictive performance.} Furthermore, the computational cost of the developed models were rarely considered in these studies, which could limit their application in settings with limited resources. In addition, external validation was \hl{reported in only two studies} \cite{Cui2024Prediction, abbas2024Explainable}\hl{, and model calibration and DCA were assessed in a few studies} \cite{Cui2024Prediction, Georgiev2024Predicting, Li2023Machine, Twait2023Dementia}\hl{. These methodological limitations could reduce the clinical utility and the reliability of the models developed in those studies for the screening of CI and dementia in real clinical settings. In contrast, the proposed framework addresses these gaps by resolving the class overlap, ensuring reliable predictions through the assessment of model calibration and decision curve analysis, and selecting models based on the performance and computational cost, while the models are further validated on an external dataset.}

\hl{The performance of the models were also evaluated against studies that used clinical data and the accuracy of the underlying assessment tools. As shown in the comparison above, despite using accessible data only, the proposed models achieved performance levels comparable to the models trained on routine clinical records} \cite{Richter2025Cognitive, Akter2025Using, Sharma2024Leveraging}\hl{. For the cognitive status tasks, the discrimination (AUC of 0.844 and 0.828) is similar to the accuracy reported for the FAQ when used alone in the same NACC data (AUC of 0.81 to 0.99)} \cite{Gonzalez2022Comprehensive, Teng2010Utility}\hl{. Furthermore, it is consistent with broader IADL accuracy across various populations} \cite{Yemm2021Instrumental, CastillaRilo2007Instrumental}\hl{. For the prediction tasks, the AUC (0.772 and 0.759) is close to the clinically validated CAIDE (Cardiovascular Risk Factors, Aging, and Incidence of Dementia) risk score (0.77)} \cite{Kivipelto2006Risk, James2021Performance}. \hl{These results indicate that the reported metrics are clinically relevant and not just artificial results produced by the models.}
